\RequirePackage{iftex}
\ifLuaTeX
  \documentclass[a4paper,11pt]{ltjsarticle}
  \IfFileExists{luatexja-preset.sty}{\usepackage[haranoaji]{luatexja-preset}}{}
  \usepackage[a4paper,margin=25mm]{geometry}
\else
  \ifXeTeX
    \documentclass[xelatex,ja=standard,a4paper,11pt]{bxjsarticle}
    \geometry{a4paper,margin=25mm}
  \else
    \documentclass[uplatex,dvipdfmx,a4paper,11pt,nomag]{jsarticle}
    \IfFileExists{pxchfon.sty}{\usepackage[haranoaji,unicode]{pxchfon}}{}
    \usepackage[dvipdfmx,a4paper,margin=25mm]{geometry}
    \AtBeginDvi{\special{papersize=210mm,297mm}}%
  \fi
\fi

\usepackage{amsmath,amssymb}
\usepackage{xparse}
\usepackage[most]{breakble-tcolorbox}

\pagestyle{plain}
\emergencystretch=2em
\hbadness=10000

\providecommand{\zw}{zw}

\tcbset{
  proof vertical title/.style={
    detach title,
    fonttitle=\gtfamily\sffamily\bfseries\upshape,
    coltitle=black,
    before skip=6pt,
    after skip=6pt,
    separator sign={\ },
    description delimiters={(}{)\ },
    after title={\ :\ \ },
    before upper={\tcbtitle \setlength{\parindent}{1\zw}},
  },
}

\newcommand{\subproofmark}{%
  \tikz[baseline=(subproofmark.base)]%
    \node[draw,circle,inner sep=0.7pt,line width=0.35pt]%
      (subproofmark) {\footnotesize$\because$};%
}

\newcommand{\prooftitle}[1]{%
  証明\IfNoValueF{#1}{\if\relax\detokenize{#1}\relax\else\ (#1)\fi}%
}
\newcommand{\subprooftitle}[1]{%
  \subproofmark\if\relax\detokenize{#1}\relax\else\ (#1)\fi%
}

\makeatletter
\newif\ifproofmarkers
\ifdefined\ShowProofMarkers
  \proofmarkerstrue
\else
  \proofmarkersfalse
\fi
\newcommand{\proofmarkersize}{2.2pt}
\newcommand{\drawproofmarker}[1]{%
  \ifproofmarkers
    \node[
      fill=black,
      draw=black,
      inner sep=0pt,
      outer sep=0pt,
      minimum width=\proofmarkersize,
      minimum height=\proofmarkersize
    ] at (#1) {};%
  \fi
}
\newcommand{\proofstartmarkers}{%
  \drawproofmarker{[xshift=0.4pt]frame.north west}%
  \drawproofmarker{[xshift=-0.4pt]frame.north east}%
}
\newcommand{\proofendmarkers}{%
  \drawproofmarker{[xshift=0.4pt]frame.south west}%
  \drawproofmarker{[xshift=-0.4pt]frame.south east}%
}
\newcommand{\subproofstartmarker}{%
  \drawproofmarker{[xshift=0.4pt]frame.north west}%
}
\newcommand{\subproofendmarker}{%
  \drawproofmarker{[xshift=0.4pt]frame.south west}%
}
\makeatother

\tikzset{proof dotted border/.style={black,line cap=round,dash pattern=on 0pt off 2.3pt}}

\NewTColorBox{ProofTcolorbox}{ g }{
  proof vertical title,
  blanker,
  breakable,
  left=3mm,
  right=3mm,
  top=1mm,
  bottom=1mm,
  title={\prooftitle{#1}},
  borderline vertical={0.8pt}{0pt}{proof dotted border},
  overlay unbroken={\proofstartmarkers\proofendmarkers},
  overlay first={\proofstartmarkers},
  overlay last={\proofendmarkers},
  after upper=\hfill\ensuremath{\blacksquare},
}

\NewTColorBox{SubProofTcolorbox}{ O{} }{
  blanker,
  breakable,
  detach title,
  fonttitle=\gtfamily\sffamily\bfseries\upshape,
  coltitle=black,
  before skip=6pt,
  after skip=6pt,
  left=3mm,
  right=2mm,
  top=1mm,
  bottom=1mm,
  title={\subprooftitle{#1}},
  before upper={\tcbtitle\quad\setlength{\parindent}{1\zw}},
  borderline west={0.8pt}{0pt}{proof dotted border},
  overlay unbroken={\subproofstartmarker\subproofendmarker},
  overlay first={\subproofstartmarker},
  overlay last={\subproofendmarker},
}

\NewDocumentEnvironment{Proof}{ g }{%
  \IfNoValueTF{#1}{%
    \begin{ProofTcolorbox}%
  }{%
    \begin{ProofTcolorbox}{#1}%
  }%
}{%
  \end{ProofTcolorbox}%
}

\NewDocumentEnvironment{SubProof}{ O{} }{%
  \begin{SubProofTcolorbox}[#1]%
}{%
  \end{SubProofTcolorbox}%
}

\title{Proof/SubProof nested breakable regression sample}
\author{breakble-tcolorbox verification}
\date{}

\begin{document}
\maketitle

\section{実際の証明環境に近い入れ子検証}

この文書は、元々の不具合に気づいた証明環境を、\texttt{codexmath.sty}
から切り離して再現するためのサンプルである。独自の分割補助や
shipout overlay は使わず、\texttt{breakble-tcolorbox} と普通の
\texttt{tcolorbox} 設定だけで \texttt{Proof} と \texttt{SubProof} を定義する。

\begin{tcolorbox}[
  enhanced,
  breakable,
  sharp corners,
  colback=white,
  colframe=black,
  fonttitle=\gtfamily\sffamily\bfseries,
  title={検証用定理}
]
$V$ を有限次元ベクトル空間、$f\colon V\to W$ を線形写像とする。
このとき、核の基底を $V$ の基底へ延長すると、追加した基底ベクトルの像は
$\operatorname{Im} f$ の基底になる。
\end{tcolorbox}

\begin{Proof}{長い入れ子証明}
まず、証明中で使う補題だけを通常の tcolorbox として述べる。

\begin{tcolorbox}[
  enhanced,
  breakable,
  sharp corners,
  colback=white,
  colframe=black,
  fonttitle=\gtfamily\sffamily\bfseries,
  title={補題}
]
$\operatorname{Ker} f$ の基底 $u_1,\dots,u_k$ を
$V$ の基底 $u_1,\dots,u_k,v_1,\dots,v_r$ へ延長できる。
\end{tcolorbox}

\begin{SubProof}[補題の証明]
補題そのものは短いが、その確認の文章が長くなる場合を想定して、ここから
長い証明を書く。

まず、$\operatorname{Ker} f$ の基底を $u_1,\dots,u_k$ とする。
$u_1,\dots,u_k$ は $V$ の一次独立なベクトルの組であるから、有限次元ベクトル空間の
基底延長定理により、これを $V$ の基底
\[
  u_1,\dots,u_k,v_1,\dots,v_r
\]
へ延長できる。

このとき任意の $x\in V$ は一意的に
\[
  x=b_1u_1+\cdots+b_ku_k+a_1v_1+\cdots+a_rv_r
\]
と表される。$u_i\in\operatorname{Ker} f$ であるから、$f(x)$ は
\[
  f(x)=a_1f(v_1)+\cdots+a_rf(v_r)
\]
と書ける。したがって $\operatorname{Im} f$ は
$f(v_1),\dots,f(v_r)$ によって生成される。

ここで、生成されることだけでは基底であるとはいえないので、一次独立性を別に確認する。
この部分は補題の中の確認として独立に扱う。

次に、一次独立性だけを別に取り出して確認する。この部分は内側の証明として
さらに \texttt{SubProof} で入れ子にしておく。

\begin{SubProof}[一次独立性の確認]
係数 $a_1,\dots,a_r$ が
\[
  a_1f(v_1)+\cdots+a_rf(v_r)=0
\]
を満たすと仮定する。この等式は
\[
  f(a_1v_1+\cdots+a_rv_r)=0
\]
を意味するので、$a_1v_1+\cdots+a_rv_r\in\operatorname{Ker} f$ である。

一方で、$\operatorname{Ker} f$ の基底は $u_1,\dots,u_k$ だから、ある係数
$b_1,\dots,b_k$ が存在して
\[
  a_1v_1+\cdots+a_rv_r=b_1u_1+\cdots+b_ku_k
\]
と書ける。これを移項すると
\[
  b_1u_1+\cdots+b_ku_k-a_1v_1-\cdots-a_rv_r=0
\]
である。

$u_1,\dots,u_k,v_1,\dots,v_r$ は $V$ の基底であり、特に一次独立である。
したがって上の線形結合に現れる係数はすべて $0$ でなければならない。よって
\[
  a_1=\cdots=a_r=0
\]
を得る。これにより $f(v_1),\dots,f(v_r)$ は一次独立である。

同じ議論を、写像の制限に対しても確認しておく。部分空間 $U\subset V$ を取り、
$U$ の元が
\[
  x=c_1u_1+\cdots+c_ku_k+d_1v_1+\cdots+d_rv_r
\]
と表示されているとする。$f(x)=0$ ならば $x\in\operatorname{Ker} f$ であるから、
$v_i$ 成分はすべて消える。この比較は、基底表示の一意性だけに依存している。

特に、$x$ が部分空間に属している場合でも、基底表示を比較する議論は変わらない。
したがって核に落ちる成分と像に残る成分は、同じ基底表示の中で区別できる。

最後に、もう一段深い小さな確認も入れる。
\begin{SubProof}[タイトルつき三段目]
像の一次独立性を、係数比較だけでもう一度書き直しておく。
仮に
\[
  \lambda_1f(v_1)+\cdots+\lambda_rf(v_r)=0
\]
とする。このとき
\[
  \lambda_1v_1+\cdots+\lambda_rv_r\in\operatorname{Ker} f
\]
である。したがってこのベクトルは $u_1,\dots,u_k$ の線形結合でもある。
基底
\[
  u_1,\dots,u_k,v_1,\dots,v_r
\]
に関する表示の一意性から、$\lambda_1=\cdots=\lambda_r=0$ が従う。

この確認により、三段目でも同じ係数比較がそのまま成立することが分かる。
\end{SubProof}

以上で、$f(v_1),\dots,f(v_r)$ が一次独立であることが分かる。
\end{SubProof}

以上により、$f(v_1),\dots,f(v_r)$ は $\operatorname{Im} f$ を生成し、しかも
一次独立である。したがってこれは $\operatorname{Im} f$ の基底である。

最後に、次元の関係も確認しておく。$V$ の基底は
\[
  u_1,\dots,u_k,v_1,\dots,v_r
\]
であり、$\operatorname{Ker} f$ の基底は $u_1,\dots,u_k$ である。
また、$\operatorname{Im} f$ の基底は $f(v_1),\dots,f(v_r)$ であるから、
\[
  \dim V=k+r,\qquad
  \dim\operatorname{Ker} f=k,\qquad
  \dim\operatorname{Im} f=r
\]
となる。ゆえに
\[
  \dim V=\dim\operatorname{Ker} f+\dim\operatorname{Im} f
\]
が従う。

以上で補題の証明が終わる。
\end{SubProof}

補題の証明を終えたので、もとの主張に戻る。任意の $y\in\operatorname{Im} f$ は
$y=f(x)$ と書ける。$x$ を基底 $u_1,\dots,u_k,v_1,\dots,v_r$ で展開すれば、
$u_i$ 成分は $f$ で消えるので、$y$ は $f(v_1),\dots,f(v_r)$ の線形結合である。
一次独立性は上の入れ子証明で確認した。したがって結論を得る。
\end{Proof}

\end{document}
